% Hinweise zum Gebrauch des Hefts (vorne im Heft).
{\fontsize{10.5}{14.5}\selectfont\setlength\parskip{2.5mm}
Das Heft zeigt jedes Programm in zwei Ansichten, beide aus denselben Daten gesetzt:

Vorne in jedem Teil steht eine Übersicht der Grade mit ihren Gruppen und Wartezeiten.

\textbf{Ein Grad je Seite.} Alles, was zu einem Grad gehört, nach Gruppen und Angriffen
geordnet. Gedacht zum Lernen und Mitnehmen.

\textbf{Matrix.} Angriff mal Grad im Querformat, wie die bisherigen Prüfungsordnungen.
Gedacht für Prüferinnen und Prüfer und für Stichproben über mehrere Grade.

\textbf{Schreibweise.} Namen in Umschrift; \textit{kursiv} steht die Ausführung oder Variante.
Längenzeichen (ô, û) wie in der jeweiligen Vorlage: Aikido ohne, Daitoryu und Shinki Toho mit.

\textbf{Japanisch.} Aikido: aus den Technik-Listen unseres Sensei. Shinki Toho: Grade und
Kamae von shinkiryu.de, Formen 1–21 von den Titelkarten der Prüfungsvideos.
Daitoryu: Die Vorlagen enthalten keine Schriftzeichen; die Angaben in \textcolor{usuzumiv}{hellem Grau}
sind ein Vorschlag und müssen noch von einem Lehrer bestätigt werden.

\textbf{Vorlagen.} Aikido Shinki Rengo, Prüfungsordnung, Stand 01.01.2022.
Daitoryu Aiki Jujutsu Bokuyokan Europa, Prüfungsordnung Kyu- und Dan-Grade, 01.04.2008.
Itto-den Shinki Toho: Grade und Prüfungsberechtigung, Kamae-Tabelle, Suburi,
Ô-dachi-Tabelle (Stand 25.08.2010) und Ô-dachi-Kommentare.
\ifentwurf

\textbf{Entwurf.} Rote Marken wie \frage{A-K-02} verweisen auf eine offene Frage in
\texttt{KORREKTUREN.md}. Im fertigen Heft fehlen sie.\fi
\par}
